\documentclass[conference]{IEEEtran}
\usepackage[utf8]{inputenc}
\usepackage{graphicx}
\usepackage{cite}
\usepackage{amsmath,amssymb,amsfonts}
\usepackage{textcomp}
\usepackage{xcolor}
\usepackage{hyperref}

\begin{document}

\title{IELTS Mock Examiner: A Full-Stack AI-Assisted Platform for Writing and Speaking Practice}

\author{
    \IEEEauthorblockN{Azlarkhon Khasakhojaev}
    \IEEEauthorblockA{Central Asian University\\Engineering School\\Department of Computer Science\\Student ID: 220263}
    \and
    \IEEEauthorblockN{Javlonbek Rustamov}
    \IEEEauthorblockA{Central Asian University\\Engineering School\\Department of Computer Science\\Student ID: 220237}
    \and
    \IEEEauthorblockN{Supervisor: [To Be Confirmed]}
    \IEEEauthorblockA{Central Asian University\\Engineering School\\Department of Computer Science\\Email: [To Be Confirmed]}
}

\maketitle

\begin{abstract}
This project presents IELTS Mock Examiner, a full-stack web platform that helps learners practice IELTS Writing and Speaking tasks with AI-assisted evaluation and interaction. The system combines a Django backend API with a React frontend and supports both standard REST workflows and real-time speaking sessions. Writing tasks are evaluated through a Gemini-based service, while speaking sessions use a streaming flow that includes voice activity detection, speech-to-text, and text-to-speech components. The backend organizes functionality into modular apps for authentication, writing, speaking, and core health/configuration services. The deployment architecture uses Docker Compose to run backend and frontend services together, with health checks and environment-driven configuration for reproducibility. The implementation focuses on practical usability, modular design, and integration of AI services in an educational setting. The resulting platform provides a foundation for structured IELTS practice and future extension toward richer analytics and adaptive feedback.
\end{abstract}

\begin{IEEEkeywords}
IELTS, AI in Education, Django, React, Speech Processing, Writing Evaluation.
\end{IEEEkeywords}

\section{Introduction}
IELTS preparation requires repeated writing and speaking practice with timely and structured feedback. Manual evaluation is time-consuming and often unavailable to many students on demand. This project addresses that gap by building a web platform that combines automated writing feedback and interactive speaking practice in one system.

The developed platform, IELTS Mock Examiner, allows users to authenticate, submit essays for task-based evaluation, and complete speaking practice using audio interaction. The system is designed to be practical for student use and maintainable for iterative academic development.

\subsection{Problem Statement}
Students preparing for IELTS often lack consistent access to immediate, structured, and criterion-oriented feedback for Writing and Speaking. Existing practice workflows are fragmented across tools and are not always aligned with an integrated learning process.

\subsection{Objectives}
The main objective is to design and implement a full-stack IELTS practice platform with AI-assisted modules. Specific objectives are:
\begin{itemize}
    \item Develop a modular backend API for authentication, writing, and speaking workflows.
    \item Implement a writing evaluation pipeline using a configurable Gemini model.
    \item Implement a speaking pipeline using VAD, STT, and TTS components.
    \item Build a responsive frontend for user interaction with writing and speaking modules.
    \item Containerize services with Docker Compose for reproducible local deployment.
\end{itemize}

\section{Methods}
\subsection{System Architecture}
The system follows a client-server architecture. The frontend is implemented with React, Vite, and TypeScript. The backend is implemented with Django, Django REST Framework, and Django Channels. The backend exposes endpoints under `/api/`, including modules for core, authentication, writing, and speaking.

Persistent storage is configured with SQLite in the current implementation through environment variable `SQLITE_PATH`. Media and cache directories are mounted in Docker for runtime assets and model caches.

\subsection{Implementation Approach}
Development uses modular application separation in the backend:
\begin{itemize}
    \item `accounts`: authentication endpoints and token-based access.
    \item `writing`: writing tasks and evaluation flow.
    \item `speaking_app`: speaking session logic and realtime interaction.
    \item `ai_services`: integrations for Gemini, STT, TTS, and VAD utilities.
    \item `core`: health and base API functionality.
\end{itemize}

The speaking flow relies on audio segmentation with configurable VAD parameters and time limits. The writing flow uses prompt templates and AI service calls to provide structured feedback.

\subsection{Technology Stack}
\textbf{Backend:} Python, Django, Django REST Framework, Django Channels.\\
\textbf{Frontend:} React, Vite, TypeScript, Tailwind CSS, shadcn/ui components.\\
\textbf{AI Services:} Gemini model configuration, Faster-Whisper STT model path configuration, Kokoro TTS model path configuration, and VAD-based segmentation.\\
\textbf{DevOps:} Docker and Docker Compose with health checks and service dependency control.

\section{Results}
\subsection{Functional Outputs}
The implemented system provides:
\begin{itemize}
    \item User authentication and session-oriented access to practice modules.
    \item Writing submission and AI-based evaluation endpoints.
    \item Speaking endpoints and realtime interaction infrastructure.
    \item Frontend pages for authentication, dashboard, writing, speaking, and administration.
\end{itemize}

\subsection{System Integration Outcomes}
Integration testing at the service level confirms that backend and frontend can run together via Docker Compose using ports 8000 and 5173, with backend health-check validation before frontend dependency startup. Environment-based configuration supports controlled model and runtime parameter adjustments.

\subsection{Current Validation Scope}
At the current project stage, validation primarily confirms feature availability and end-to-end service connectivity. The platform is ready for extended experimental evaluation such as latency benchmarking, scoring consistency analysis, and user acceptance testing.

\section{Discussion}
The project demonstrates that an integrated IELTS practice system can be built with a modular web architecture and configurable AI services. A key strength is separation of concerns across backend apps and frontend pages, which simplifies maintenance and extension.

The current limitations include the absence of formal large-scale user testing and benchmark reporting in this version of the report. In addition, model output quality may vary depending on prompt design and runtime configuration.

Despite these limitations, the system establishes a solid technical baseline for incremental improvements in evaluation accuracy, speaking interaction smoothness, and pedagogical feedback quality.

\section{Conclusion}
IELTS Mock Examiner successfully implements a full-stack educational platform for writing and speaking practice with AI-assisted processing. The backend, frontend, and containerized deployment are functionally integrated and aligned with the project objectives.

Future work includes rubric-level quantitative evaluation, richer analytics dashboards, expanded speaking session orchestration, and comparative testing of model configurations for quality and performance.

\begin{thebibliography}{9}
\bibitem{django}
Django Software Foundation, "Django Documentation," https://docs.djangoproject.com/.

\bibitem{drf}
Encode OSS, "Django REST Framework Documentation," https://www.django-rest-framework.org/.

\bibitem{react}
Meta, "React Documentation," https://react.dev/.

\bibitem{vite}
Vite Team, "Vite Documentation," https://vite.dev/.

\bibitem{channels}
Django Channels Contributors, "Channels Documentation," https://channels.readthedocs.io/.
\end{thebibliography}

\end{document}
